\chapter{Analytic maps}
\label{chap:analytic}

This chapter is a plan, and its statements become formal one by one. We follow the treatment of
Glöckner and Neeb \cite[Ch.~2]{GN26}, which builds on Bochnak and Siciak \cite{BS71a, BS71b} and
on Glöckner \cite[§2]{Glo02}. The point of this approach is that \emph{no completeness of the
target space is needed} for the basic theory. Bochnak and Siciak assume throughout that the target
space $F$ is sequentially complete; Glöckner removes this hypothesis by working in the completion
$\widetilde F$ and coming back to $F$ through the derivatives, which lie in $F$
\cite[Prop.~2.4, Lemma~2.5]{Glo02}.

\section{Continuous homogeneous polynomials}

\begin{definition}[Continuous homogeneous polynomials]
  \label{def:hom-poly}
  \lean{LCAnalysis.IsContHomPoly}
  \leanok
  Let $E$, $F$ be topological vector spaces over $\mathbb{K} \in \{\mathbb{R}, \mathbb{C}\}$ and
  $k \in \mathbb{N}_0$. A map $p \colon E \to F$ is a \emph{continuous homogeneous polynomial of
  degree $k$} if $p(x) = A(x, \dots, x)$ for some continuous $k$-linear map
  $A \colon E^k \to F$.
\end{definition}

Replacing $A$ by its symmetrization we may take $A$ symmetric, and then $A$ is determined by $p$:

\begin{theorem}[Polarization; Mazur--Orlicz {\cite[Theorem~A]{BS71a}}]
  \label{thm:polarization}
  \lean{MultilinearMap.polarization, MultilinearMap.continuous_of_isSymm_of_continuous_diag,
    LCAnalysis.isContHomPoly_iff_continuous}
  \leanok
  If $A$ is symmetric and $k$-linear, then for $x_1, \dots, x_k \in E$
  \[
    \sum_{\varepsilon \in \{\pm 1\}^k} \varepsilon_1 \cdots \varepsilon_k \,
    A\Bigl(\sum_j \varepsilon_j x_j, \dots, \sum_j \varepsilon_j x_j\Bigr)
    = 2^k k! \, A(x_1, \dots, x_k).
  \]
  In particular, if $x \mapsto A(x, \dots, x)$ is continuous, then $A$ is continuous.
\end{theorem}

\begin{proof}
  \leanok
  Expand $A(\sum_j \varepsilon_j x_j, \dots)$ by multilinearity over all maps
  $r \colon \{1, \dots, k\} \to \{1, \dots, k\}$. The sign sum
  $\sum_\varepsilon \prod_j \varepsilon_j^{1 + |r^{-1}(j)|}$ is $2^k$ if $r$ is a bijection and
  $0$ otherwise, and $A(x_{r(1)}, \dots, x_{r(k)}) = A(x_1, \dots, x_k)$ for bijective $r$ by
  symmetry. The continuity statement follows because the left hand side is a finite combination of
  values of $x \mapsto A(x, \dots, x)$. No completeness, local convexity or Hausdorff property is
  used.
\end{proof}

\section{One complex variable}

Let $F$ be a complex locally convex space, not necessarily complete, and $U \subseteq \mathbb{C}$
open. A map $f \colon U \to F$ is \emph{complex analytic} if around every point it is a power
series $f(z_0 + z) = \sum_n z^n a_n$ with $a_n \in F$, converging pointwise on a disc
\cite[Rem.~2.1.3]{GN26}.

\begin{theorem}[Power series in incomplete spaces]
  \label{thm:one-var-power-series}
  \lean{LCAnalysis.HasPowerSeriesOnDisc.continuousOn, LCAnalysis.CAnalyticOn.cAnalyticOn_cDeriv,
    LCAnalysis.HasPowerSeriesOnDisc.hasWeakIntegral, LCAnalysis.HasPowerSeriesOnDisc.seminorm_coeff_le,
    LCAnalysis.CAnalyticOn.hasPowerSeriesOnDisc}
  \leanok
  A complex analytic $f \colon U \to F$ is continuous, its complex derivative exists and is
  complex analytic. The Cauchy integrals
  $\int_0^{2\pi} e^{-int} f(z_0 + \rho e^{it})\, dt = 2\pi \rho^n a_n$ exist as weak integrals in
  $F$, the Cauchy estimates $q(a_n) \le \rho^{-n} \sup_{|\zeta| = \rho} q(f(z_0 + \zeta))$ hold
  for every continuous seminorm $q$, and the Taylor series of $f$ at $z_0$ converges to $f$ in $F$
  on every disc whose closure lies in $U$ (\cite[Cor.~2.1.8, Lemma~2.1.10]{GN26}).
\end{theorem}

\begin{proof}
  \leanok
  No completion is used. Every sum and integral that appears has its value already in $F$ (a value
  of $f$, a coefficient, a derivative), and convergence is checked seminorm by seminorm against a
  geometric majorant. Scalar estimates are transferred to $F$ by the Hahn--Banach theorem: if
  $|\ell(v)| \le B$ for every continuous linear $\ell$ with $|\ell| \le q$, then $q(v) \le B$.
\end{proof}

\begin{theorem}[Complex analytic $=$ real smooth with complex linear differential]
  \label{thm:one-var-iff}
  \lean{LCAnalysis.cAnalyticOn_iff_smoothMBOn}
  \leanok
  \uses{thm:one-var-power-series}
  $f \colon U \to F$ is complex analytic if and only if it is $C^\infty$ in the real
  Michal--Bastiani sense and $\dd f(z)(w) = w\, \dd f(z)(1)$ for all $z \in U$, $w \in \mathbb{C}$
  \cite[Cor.~2.1.9]{GN26}.
\end{theorem}

\begin{remark}
  Glöckner and Neeb prove the one variable characterizations with $C^1$ in place of $C^\infty$
  under Mackey completeness \cite[Prop.~2.1.6]{GN26}, and Corollary~2.1.9 without completeness
  through the completion $\widetilde F$. Here the smooth version is proved without completeness
  and without the completion: smoothness provides all the Taylor coefficients
  $\frac{1}{n!}\dd^{(n)} f(z)(1, \dots, 1)$ in $F$. For a single power series on a disc, the
  derivative at a point other than the center is only known to exist when the differentiated
  series converges in $F$; whether this can fail in an incomplete space we have not settled.
\end{remark}

\section{Complex analytic maps}

\begin{definition}[Complex analytic maps {\cite[Def.~2.1.2]{GN26}, \cite[Def.~5.6]{BS71b}}]
  \label{def:complex-analytic}
  Let $E$, $F$ be complex locally convex spaces and $U \subseteq E$ open. A map
  $f \colon U \to F$ is \emph{complex analytic} if it is continuous and for every $x \in U$ there
  are continuous homogeneous polynomials $p_k$ of degree $k$ such that
  $f(x + y) = \sum_{k \ge 0} p_k(y)$, as a pointwise limit, for all $y$ in a $0$-neighborhood $Y$
  with $x + Y \subseteq U$.
\end{definition}

\begin{definition}[Gâteaux analytic maps {\cite[Def.~2.1.11]{GN26}}]
  \label{def:gateaux-analytic}
  $f$ is \emph{Gâteaux analytic} if $z \mapsto f(x + z y)$ is complex analytic (locally a power
  series with coefficients in $F$) on $\{z \in \mathbb{C} : x + zy \in U\}$, for all $x \in U$,
  $y \in E$.
\end{definition}

\begin{theorem}[Characterizations of complex analyticity {\cite[Thm.~2.1.12]{GN26}}]
  \label{thm:complex-analytic-iff}
  For $f \colon U \to F$ the following are equivalent:
  \begin{enumerate}
    \item $f$ is $C^\infty$ (in the real Michal--Bastiani sense) and $\dd f(x) \colon E \to F$ is
      complex linear for each $x \in U$;
    \item $f$ is continuous and Gâteaux analytic;
    \item $f$ is complex analytic.
  \end{enumerate}
  If $F$ is Mackey complete, they are also equivalent to: $f$ is $C^1$ and $\dd f(x)$ is complex
  linear for each $x \in U$.
\end{theorem}

\begin{remark}
  A word of caution on the hypotheses. No completeness is needed for the equivalence of (1)--(3):
  in the proof of (2) $\Rightarrow$ (1), the Cauchy integral
  $\dd f(x, y) = \frac{1}{2\pi i} \oint_{|\zeta| = 1} f(x + \zeta y)\, \zeta^{-2}\, d\zeta$ is a
  weak integral which exists in $F$, because along the line it is a coefficient of a power series
  with coefficients in $F$. Completeness only enters when one starts from a $C^1$ map, whose
  higher derivatives are a priori only in the completion. Since (1) is stated with the real
  calculus, we do not need a separate complex Michal--Bastiani calculus: a complex $C^k$ map is
  the same as a real $C^k$ map with complex linear differentials.
\end{remark}

\section{Real analytic maps}

\begin{definition}[Complexification {\cite[Def.~2.2.1]{GN26}}]
  \label{def:complexification}
  For a real topological vector space $E$, $E_{\mathbb{C}} := E \times E$ with the product
  topology and $(x + iy)(u, v) := (xu - yv, xv + yu)$. It is locally convex if $E$ is, and
  $E = E \times \{0\}$.
\end{definition}

\begin{definition}[Real analytic maps; Milnor {\cite{Mil84}}, {\cite[Def.~2.3]{Glo02}},
  {\cite[Def.~2.2.2]{GN26}}]
  \label{def:real-analytic}
  Let $E$, $F$ be real locally convex spaces and $U \subseteq E$ open. A map $f \colon U \to F$ is
  \emph{real analytic} if it extends to a complex analytic map $V \to F_{\mathbb{C}}$ on an open
  neighborhood $V$ of $U$ in $E_{\mathbb{C}}$.
\end{definition}

\begin{theorem}[Basic properties {\cite[Lemmas~2.2.3, 2.2.5, 2.2.6, Prop.~2.2.4]{GN26}}]
  \label{thm:real-analytic-basic}
  Real analyticity is a local property; compositions of real analytic maps are real analytic; real
  analytic maps are smooth; and a real analytic map is locally the sum of its Taylor series
  $\sum_k \frac{1}{k!} \dd^{(k)} f(x)(y, \dots, y)$.
\end{theorem}

\begin{remark}[Two notions of real analyticity]
  Bochnak and Siciak call a real map analytic when it is continuous and locally a pointwise sum of
  continuous homogeneous polynomials \cite[Def.~5.6]{BS71b}. Milnor's definition above is
  stronger. The two agree when $E_{\mathbb{C}}$ is a Baire space, e.g.\ for Fréchet spaces
  \cite[Thm.~7.1]{BS71b}, \cite[2.9]{Glo02}. For the weaker notion, closedness under composition
  is only known when the spaces are Baire \cite[Thm.~7.3]{BS71b}; this is why the theory of
  analytic manifolds and Lie groups uses Milnor's definition.
\end{remark}

\section{Comparison with Mathlib on Banach spaces}

On Banach spaces, Mathlib's \texttt{AnalyticOnNhd} (convergent power series in norm) should be
compared with the definitions above, as Theorem~\ref{thm:mb-to-frechet} compares the Michal--Bastiani and
Fréchet calculi. This is planned once the locally convex theory is in place.
